% 第7批：《大学物理简明教程》第7章静电场解答题7.3.15—7.3.20。
% 原书/PDF第188页，用户扫描件校对；独立推导及限制条件。
\subsection{解答题：导体静电平衡、电介质与场能}

\sourceproblem{7.3.15}{188}{188}
\problemtext 半径分别为$R_1<R_2$的两同心薄金属球壳，起初内球壳净带电$+q$、
外球壳不带电。
(1) 求外壳内、外表面电荷分布以及外壳电势；
(2) 将外壳接地后再断开地线，求此时外壳电荷分布及电势。
\workings
球对称性及静电平衡使内壳电荷$+q$分布在其外表面。
(1) 在外壳金属内取高斯面，包围电荷为零，故外壳内表面总电荷
$\boxed{Q_{2,\mathrm{in}}=-q}$。
因为外壳原净电荷为零，外壳外表面
$\boxed{Q_{2,\mathrm{out}}=+q}$，
均匀分布于同心球面。
外球壳电势（无穷远为零）
\[
\boxed{\varphi_2=\frac{q}{4\pi\varepsilon_0R_2}},
\]
因为在$r\ge R_2$处系统的总电荷为$q$。
(2) 外壳接地时依然必须有内表面电荷$-q$；
地线使自由电荷流出，外壳外表面净电荷变成零。
因此
\[
\boxed{Q_{2,\mathrm{in}}=-q,\quad
Q_{2,\mathrm{out}}=0,\quad
Q_{2,\mathrm{net}}=-q},\qquad
\boxed{\varphi_2=0}.
\]
\crosscheck 接地后系统在最外部的净包围电荷为零，
故外部静电场与电势均为零。
断开接地线后这个静电平衡解保持（无其他扰动）。

\sourceproblem{7.3.16}{188}{188}
\problemtext 半径为$R_1$的带电金属球总电荷$Q$，
外围包有从$r=R_1$到$r=R_2$的均匀线性电介质球壳，
相对介电常数$\varepsilon_r$，球壳外为真空。
求各区域电场强度、电势以及金属球电势，
以无穷远为零电势。
\workings
对任意半径$r>R_1$的球面，电位移矢量的通量只由自由电荷$Q$决定，
\[
D_r(r)=\frac{Q}{4\pi r^2}.
\]
分段得到
\[
\boxed{\bm E(r)=
\begin{cases}
\bm0,&0\le r<R_1,\\
\dfrac{Q}{4\pi\varepsilon_0\varepsilon_r r^2}\bm e_r,
  &R_1<r<R_2,\\
\dfrac{Q}{4\pi\varepsilon_0 r^2}\bm e_r,&r>R_2.
\end{cases}}
\]
使用$\varphi(r)=\int_r^\infty E_r(s)\,ds$，
外真空区
\[
\boxed{\varphi(r)=\frac{Q}{4\pi\varepsilon_0r}},\quad r\ge R_2.
\]
介质内部
\[
\boxed{\varphi(r)=\frac{Q}{4\pi\varepsilon_0}
\left[\frac1{R_2}+
\frac1{\varepsilon_r}\left(\frac1r-\frac1{R_2}\right)\right]},
\quad R_1\le r\le R_2.
\]
金属球内电势处处等于表面值
\[
\boxed{\varphi_{\mathrm{metal}}=
\frac{Q}{4\pi\varepsilon_0}
\left[\frac1{\varepsilon_rR_1}
+\left(1-\frac1{\varepsilon_r}\right)\frac1{R_2}\right]}.
\]
\crosscheck $\varepsilon_r=1$时退化为真空带电金属球；
在$r=R_2$的两侧电势连续，且$D_r$连续，
$E_r$则因材料介电常数改变而出现跳变。

\sourceproblem{7.3.17}{188}{188}
\problemtext 图7.3.17中，一个中性空心导体$A$的同一空腔内有两个相互绝缘的导体$B,C$。
$B$带净正电，$A$和$C$不带净电。
比较三导体静电平衡时的电势$U_A,U_B,U_C$。
\workings
将空腔外导体$A$作为参考电势$U_A$。
$B$带净正电，其电势相对于边界$A$为正；
带正电的$B$在腔内产生电场。
处于相同腔内的中性浮置导体$C$会极化，
其电势在静电平衡时处处为常数，取值由零净电荷条件决定。
可把$C$的电势看成周围固定电势边界的正权平均；
因此在典型的连通、无其他带电体的构型中
\[
\boxed{U_B>U_C>U_A}.
\]
\crosscheck 上述严格比较依赖$B,C$均独立、互不接触、彼此不屏蔽，
且腔内只有$B$携带非零自由净电荷。
若另施加外部电势边界或其他腔内带电体，则需重新求解静电边值问题。

\sourceproblem{7.3.18}{188}{188}
\problemtext 真空中均匀带电实心球体半径$R$、总电荷$q$，
按电场能密度$w_e=\tfrac12\varepsilon_0E^2$计算其全部静电场能量。
\workings
均匀体电荷密度为$\rho=3q/(4\pi R^3)$。
由高斯定理得到
\[
E(r)=
\begin{cases}
\dfrac{qr}{4\pi\varepsilon_0R^3},&0\le r<R,\\
\dfrac{q}{4\pi\varepsilon_0r^2},&r\ge R.
\end{cases}
\]
用球壳体积元$dV=4\pi r^2dr$积分：
\[
W_{\mathrm{in}}=\frac{\varepsilon_0}{2}
\int_0^R E(r)^2\,4\pi r^2dr
=\frac{q^2}{40\pi\varepsilon_0R},
\]
\[
W_{\mathrm{out}}=\frac{\varepsilon_0}{2}
\int_R^\infty E(r)^2\,4\pi r^2dr
=\frac{q^2}{8\pi\varepsilon_0R}.
\]
所以
\[
\boxed{W_e=W_{\mathrm{in}}+W_{\mathrm{out}}
=\frac{3q^2}{20\pi\varepsilon_0R}}.
\]
\crosscheck 均匀带电\emph{实体}球的静电能量是
$3k_eq^2/(5R)$；若所有电荷只分布在球面，则是$k_eq^2/(2R)$，
二者不要混淆。

\sourceproblem{7.3.19}{188}{188}
\problemtext 真空中平行板电容器正对面积
$S=20.0\,\mathrm{cm^2}$，板间距$d=2.00\,\mathrm{mm}$，
两板电势差$U=400\,\mathrm V$。
忽略边缘效应，求电容器储存的静电能量。
\workings
换算$S=2.00\times10^{-3}\,\mathrm{m^2}$、
$d=2.00\times10^{-3}$ m，
\[
C=\frac{\varepsilon_0S}{d}
=\varepsilon_0\approx8.85\times10^{-12}\,\mathrm F.
\]
因此
\[
\boxed{W_e=\frac12CU^2
=\frac12(8.85\times10^{-12})(400)^2
\approx7.08\times10^{-7}\,\mathrm J}.
\]
\crosscheck 两种等价公式
$W_e=\tfrac12 QU=\tfrac{Q^2}{2C}$也成立；
这里电压单位为伏、能量单位为焦耳。

\sourceproblem{7.3.20}{188}{188}
\problemtext 真空中的球形电容器由两个同心金属球面组成，
内、外半径分别$R_1<R_2$，所带电荷分别为$+Q,-Q$。
求该球形电容器的电容和储能。
\workings
在$R_1<r<R_2$区域，
\[
E(r)=\frac{Q}{4\pi\varepsilon_0r^2}.
\]
两板电势差为
\[
U=\int_{R_1}^{R_2}E(r)\,dr
=\frac{Q}{4\pi\varepsilon_0}
\left(\frac1{R_1}-\frac1{R_2}\right).
\]
所以
\[
\boxed{C=\frac Q U
=\frac{4\pi\varepsilon_0R_1R_2}{R_2-R_1}},
\]
\[
\boxed{W_e=\frac{Q^2}{2C}
=\frac{Q^2}{8\pi\varepsilon_0}
\left(\frac1{R_1}-\frac1{R_2}\right)}.
\]
\crosscheck $R_2\to\infty$时，电容趋于孤立导体球电容$4\pi\varepsilon_0R_1$；
当两壳间距缩小时，电容增大。
